% !TEX program = xelatex
% ============================================================
%  اتفاقية الشراكة — Partnership Agreement Template
%  نموذج اتفاقية شراكة بين طرفين أو أكثر لتأسيس وإدارة مشروع مشترك
%  Compile with: xelatex main.tex (run twice)
% ============================================================
%  Features:
%    - أطراف ومقدمات (Parties & Recitals)
%    - مواد مرقمة: الغرض، المساهمات، توزيع الأرباح، الإدارة، الواجبات، المدة، الإنهاء، حل النزاعات
%    - جدول مساهمات الشركاء (Contributions Table)
%    - جدول توزيع الأرباح والخسائر (Profit/Loss Sharing Table)
%    - كتلة توقيعات للأطراف (Signatures Block)
%    - نص عربي قانوني مع مصطلحات إنجليزية بين قوسين
% ============================================================

\documentclass[12pt,a4paper]{article}

\usepackage[a4paper,margin=2.5cm,top=2.5cm,bottom=2.5cm,headheight=16pt]{geometry}
\usepackage{graphicx}
\usepackage{array}
\usepackage{booktabs}
\usepackage{tabularx}
\usepackage{multirow}
\usepackage{enumitem}
\usepackage{float}
\usepackage{titlesec}
\usepackage{fancyhdr}
\usepackage{setspace}
\usepackage{xcolor}

\usepackage{bidi}
\usepackage[no-math]{fontspec}
\usepackage[quiet]{polyglossia}

\setmainfont[Scale=1]{NotoKufiArabic-Regular.ttf}
\newfontfamily\arabicfont[Script=Arabic,Scale=0.95]{NotoKufiArabic-Regular.ttf}
\newfontfamily\englishfont[Scale=0.95]{TeX Gyre Termes}

\setdefaultlanguage[calendar=gregorian,hijricorrection=1]{arabic}
\setotherlanguage[variant=british]{english}

\usepackage[breaklinks,hidelinks]{hyperref}
\hypersetup{pdftitle={اتفاقية الشراكة},pdfauthor={اسم الشركة}}

\titleformat{\section}{\large\bfseries}{\thesection}{0.7em}{}
\titlespacing*{\section}{0pt}{1.2ex plus 0.3ex}{0.8ex plus 0.2ex}

\pagestyle{fancy}
\fancyhf{}
\fancyhead[R]{\small\itshape اتفاقية الشراكة (Partnership Agreement)}
\fancyhead[L]{\small اسم الشركة}
\fancyfoot[C]{\small\thepage}
\renewcommand{\headrulewidth}{0.3pt}

\newcolumntype{L}[1]{>{\raggedright\arraybackslash}p{#1}}
\newcolumntype{C}[1]{>{\centering\arraybackslash}p{#1}}

\setstretch{1.3}

\begin{document}

\begin{center}
{\LARGE\bfseries اتفاقية الشراكة}\\[0.5em]
{\large (\LR{Partnership Agreement})}\\[1em]
\end{center}

\vspace{0.5em}

% TODO: استبدل التاريخ والمرجع بالقيم الفعلية
\noindent\textbf{رقم الاتفاقية:} \rule{4cm}{0.4pt} \hfill \textbf{التاريخ:} \rule{4cm}{0.4pt}

\vspace{1em}

% ============================================================
% PARTIES
% ============================================================
\noindent\textbf{الأطراف (\LR{Parties}):}

\noindent وقّع الطرفان أدناه على هذه الاتفاقية، ويُشار إليهما مجتمعين بـ«الطرفين» ومنفردين بـ«الطرف»:

\begin{enumerate}
    \item \textbf{الطرف الأول (\LR{First Party}):} \rule{6cm}{0.4pt}، الشركة المسجّلة في \rule{4cm}{0.4pt}، رقم السجل التجاري: \rule{3cm}{0.4pt}، ويمثلها في التوقيع \rule{5cm}{0.4pt} بصفته \rule{3cm}{0.4pt}.

    \item \textbf{الطرف الثاني (\LR{Second Party}):} \rule{6cm}{0.4pt}، رقم الهوية/الجواز: \rule{3cm}{0.4pt}، العنوان: \rule{5cm}{0.4pt}.
\end{enumerate}

\vspace{1em}

% ============================================================
% RECITALS
% ============================================================
\noindent\textbf{الديباجة (\LR{Recitals}):}

\begin{enumerate}
    \item يمتلك الطرف الأول خبرة وموارد في مجال \rule{5cm}{0.4pt}.
    \item يمتلك الطرف الثاني خبرة وموارد في مجال \rule{5cm}{0.4pt}.
    \item يرغب الطرفان في تأسيس شراكة (\LR{Partnership}) لإدارة مشروع مشترك وفقاً لشروط هذه الاتفاقية.
    \item وقّع الطرفان هذه الاتفاقية بمحض إرادتهما الحرّة وبكامل أهليتهما القانونية.
\end{enumerate}

\vspace{1em}

% ============================================================
% 1. PURPOSE
% ============================================================
\section{المادة الأولى: الغرض من الشراكة (\LR{Purpose})}
\label{sec:purpose}
يتمثل الغرض من هذه الشراكة في تأسيس وإدارة مشروع مشترك (\LR{Joint Venture}) في مجال \rule{6cm}{0.4pt}، وتشمل الأهداف ما يلي:
\begin{enumerate}
    \item تطوير وتسويق المنتجات/الخدمات التالية: \rule{8cm}{0.4pt}.
    \item تحقيق عوائد مالية مشتركة للطرفين.
    \item الاستفادة من خبرات وموارد الطرفين بشكل تكاملي.
\end{enumerate}

% ============================================================
% 2. CONTRIBUTIONS
% ============================================================
\section{المادة الثانية: مساهمات الشركاء (\LR{Contributions})}
\label{sec:contributions}
يلتزم كل طرف بتقديم المساهمات المحددة في الجدول أدناه خلال \rule{2cm}{0.4pt} يوماً من تاريخ توقيع الاتفاقية:

% TODO: املأ تفاصيل مساهمات كل طرف
\begin{table}[H]
\centering
\renewcommand{\arraystretch}{1.5}
\begin{tabularx}{\textwidth}{L{4cm} X X}
\toprule
\textbf{نوع المساهمة} & \textbf{الطرف الأول} & \textbf{الطرف الثاني} \\
\midrule
المساهمة النقدية (\LR{Cash}) & \rule{4cm}{0.4pt} & \rule{4cm}{0.4pt} \\
المساهمة العينية (\LR{In-kind}) & \rule{4cm}{0.4pt} & \rule{4cm}{0.4pt} \\
الخبرة والعمل (\LR{Labor}) & \rule{4cm}{0.4pt} & \rule{4cm}{0.4pt} \\
\midrule
\textbf{إجمالي القيمة} & \textbf{\rule{4cm}{0.4pt}} & \textbf{\rule{4cm}{0.4pt}} \\
\bottomrule
\end{tabularx}
\end{table}

% ============================================================
% 3. PROFIT/LOSS SHARING
% ============================================================
\section{المادة الثالثة: توزيع الأرباح والخسائر (\LR{Profit/Loss Sharing})}
\label{sec:profitloss}
\begin{enumerate}
    \item تُوزَّع الأرباح والخسائر (\LR{Profits and Losses}) بين الطرفين بالنسب المحددة في الجدول أدناه:

% TODO: حدّد نسب التوزيع
\begin{table}[H]
\centering
\renewcommand{\arraystretch}{1.5}
\begin{tabularx}{\textwidth}{L{6cm} C{3cm} C{3cm}}
\toprule
\textbf{البند} & \textbf{الطرف الأول} & \textbf{الطرف الثاني} \\
\midrule
نسبة توزيع الأرباح & \rule{2cm}{0.4pt}\% & \rule{2cm}{0.4pt}\% \\
نسبة تحمّل الخسائر & \rule{2cm}{0.4pt}\% & \rule{2cm}{0.4pt}\% \\
\bottomrule
\end{tabularx}
\end{table}

    \item تُوزَّع الأرباح \rule{2cm}{0.4pt} (شهرياً/ربعياً/سنوياً) بناءً على القوائم المالية المعتمدة.
    \item تُرحَّل الخسائر إلى السنوات المالية التالية وفق القوانين السارية.
\end{enumerate}

% ============================================================
% 4. MANAGEMENT
% ============================================================
\section{المادة الرابعة: إدارة الشراكة (\LR{Management})}
\label{sec:management}
\begin{enumerate}
    \item تُدار الشراكة بقرارات مشتركة بين الطرفين، وتتطلب القرارات الكبرى موافقة الطرفين بالكتابة.
    \item تُتخذ القرارات اليومية (\LR{Day-to-day Operations}) من قبل \rule{5cm}{0.4pt}.
    \item تُعقد اجتماعات دورية \rule{2cm}{0.4pt} (شهرياً/ربعياً) لمراجعة الأداء المالي والتشغيلي.
    \item تُحفظ سجلات (\LR{Records}) دقيقة لجميع القرارات والاجتماعات.
\end{enumerate}

% ============================================================
% 5. DUTIES
% ============================================================
\section{المادة الخامسة: واجبات الشركاء (\LR{Duties})}
\label{sec:duties}
يلتزم كل طرف بما يلي:
\begin{enumerate}
    \item تقديم المساهمات المتفق عليها في الوقت المحدد.
    \item العمل بأمانة وإخلاص لتحقيق أهداف الشراكة (\LR{Fiduciary Duty}).
    \item عدم الانخراط في أنشطة تتعارض مع مصالح الشراكة.
    \item الحفاظ على سرية معلومات الشراكة.
    \item عدم نقل حقوق أو التزامات هذه الاتفاقية لطرف ثالث دون موافقة الطرف الآخر كتابةً.
\end{enumerate}

% ============================================================
% 6. TERM
% ============================================================
\section{المادة السادسة: مدة الاتفاقية (\LR{Term})}
\label{sec:term}
\begin{enumerate}
    \item \textbf{تاريخ بدء الاتفاقية:} \rule{4cm}{0.4pt}
    \item \textbf{مدة الاتفاقية:} \rule{2cm}{0.4pt} سنوات، قابلة للتجديد بالاتفاق المتبادل.
    \item \textbf{الإشعار قبل الانتهاء:} يُشعر الطرف الآخر قبل \rule{2cm}{0.4pt} يوماً من تاريخ الانتهاء في حال عدم الرغبة في التجديد.
\end{enumerate}

% ============================================================
% 7. TERMINATION
% ============================================================
\section{المادة السابعة: إنهاء الاتفاقية (\LR{Termination})}
\label{sec:termination}
يجوز إنهاء هذه الاتفاقية في الحالات التالية:
\begin{enumerate}
    \item \textbf{انتهاء المدة:} بانتهاء المدة المتفق عليها دون تجديد.
    \item \textbf{الاتفاق المتبادل:} بموافقة الطرفين كتابةً.
    \item \textbf{الإخلال الجسيم:} عند إخلال أحد الطرفين بالتزاماته، مع منحه مهلة \rule{2cm}{0.4pt} يوماً للتصحيح.
    \item \textbf{الإفلاس أو التصفية:} عند إفلاس أحد الطرفين أو تصفية أعماله.
\end{enumerate}

\noindent عند الإنهاء، تُصفّى أصول الشراكة وتُسدد الالتزامات، ثم يُوزَّع ما تبقّى بين الطرفين بالنسب المتفق عليها.

% ============================================================
% 8. DISPUTE RESOLUTION
% ============================================================
\section{المادة الثامنة: حل النزاعات (\LR{Dispute Resolution})}
\label{sec:dispute}
\begin{enumerate}
    \item يلتزم الطرفان بحل النزاعات (\LR{Disputes}) ودياً عبر المفاوضات المباشرة أولاً.
    \item في حال عدم الوصول إلى حل خلال \rule{2cm}{0.4pt} يوماً، يُحال النزاع إلى التحكيم (\LR{Arbitration}) في \rule{4cm}{0.4pt}.
    \item تُفسَّر هذه الاتفاقية وفقاً لقوانين \rule{5cm}{0.4pt}.
\end{enumerate}

% ============================================================
% SIGNATURES
% ============================================================
\vspace{2em}
\section*{التوقيعات (\LR{Signatures})}

% TODO: املأ بيانات الموقّعين
\begin{table}[H]
\centering
\renewcommand{\arraystretch}{2.5}
\begin{tabularx}{\textwidth}{L{5cm} X X}
\toprule
& \textbf{الطرف الأول} & \textbf{الطرف الثاني} \\
\midrule
\textbf{الاسم} & --- & --- \\
\textbf{الصفة} & --- & --- \\
\textbf{التوقيع} & \rule{4cm}{0.4pt} & \rule{4cm}{0.4pt} \\
\textbf{التاريخ} & \rule{4cm}{0.4pt} & \rule{4cm}{0.4pt} \\
\bottomrule
\end{tabularx}
\end{table}

\end{document}
